% ============================================================================
%  COVER LETTER -- Radiology: Artificial Intelligence (RSNA)
%
%  THIS FILE COMPILES ON ITS OWN. Engine: pdfLaTeX. One pass is enough.
%  Only standard packages. No graphics, no bibliography.
%
%  BEFORE SENDING: fill in the two placeholders marked <<< >>> below
%  (submission date; archival DOI once minted).
% ============================================================================
\documentclass[11pt,a4paper]{article}

\usepackage[T1]{fontenc}
\usepackage[utf8]{inputenc}
\usepackage{lmodern}
\usepackage[margin=1in]{geometry}
\usepackage{microtype}
\usepackage{parskip}
\usepackage[hidelinks]{hyperref}

\pagestyle{empty}
\setlength{\emergencystretch}{2em}

\begin{document}

\noindent
\textbf{Sathvik Loke}\\
Illinois Mathematics and Science Academy\\
1500 Sullivan Road, Aurora, IL 60506, USA\\
\texttt{sloke@imsa.edu}

\vspace{1.2em}

\noindent
<<< SUBMISSION DATE >>>

\vspace{1.2em}

\noindent
The Editor\\
\emph{Radiology: Artificial Intelligence}\\
Radiological Society of North America

\vspace{1.2em}

\noindent Dear Editor,

We submit for your consideration \textbf{``What a slice-level benchmark certifies
without the pixels: a label-file audit of seven public benchmarks, and a
protocol-frozen random-sample screen of how often the check is reported''} as an
Original Research article.

\subsection*{What is already known, stated first}

The central construction in this paper --- fitting a model to slice position alone,
with no image data, and comparing it with a published slice-level number --- is
\textbf{prior art}, and the manuscript says so in its introduction rather than in a
late discussion paragraph. We would rather you learn this from us than from a
reviewer.

\begin{itemize}
\item A Kaggle notebook by the user OsciiArt, ``Baseline with no image'' (RSNA-STR
Pulmonary Embolism Detection, October 2020, gold medal), bins the label rate against
relative slice position on the training set, applies it to the test set with no
pixels, and scores 0.33 on the public leaderboard against 0.44 for a constant
predictor. That is the same construction, four years earlier.
\item Yan et al.\ (CVPR 2018), in the paper that defines the DeepLesion lesion-type
task, report a ``Baseline: Location feature'' row at 59.7\% against their full
method's 90.5\%. A location-only baseline on that benchmark is theirs.
\item Badgeley et al.\ (\emph{npj Digital Medicine} 2019;2:31) showed metadata-driven
performance in hip-fracture classification collapsing to AUC 0.52 on a
confounder-balanced test set, with scanner model predictable at AUC 1.00.
\item Zech et al.\ (\emph{PLOS Medicine} 2018) and Geirhos et al.\ (\emph{Nature
Machine Intelligence} 2020) established shortcut learning as the framing.
\item Yagis et al.\ (2021) and Tampu et al.\ (2022) measured the inflation caused by
slice-level evaluation.
\item Ong Ly et al.\ (\emph{npj Digital Medicine} 2024) is the closest competitor:
thirteen datasets, and a released estimator.
\item Maier-Hein et al.\ (\emph{Nature Communications} 2018;9:5217) already showed
that challenge rankings are unstable, and that Kendall's $\tau$ in the 0.74--0.85
range is enough for critical changes in ranking. This caps what anyone may claim
about rank inversion, and it caps what we claim.
\end{itemize}

\subsection*{What is new}

Five things, and we have tried to size each of them honestly.

\textbf{1. A patient-level result that needs no published comparator to be read.} On
the RSNA 2019 Intracranial Haemorrhage benchmark --- 18{,}938 patients, 752{,}802
slices, 21{,}744 series --- a zero-image positional baseline reaches slice-level AUROC
0.737 [0.735, 0.740] and patient-level AUROC \textbf{0.453} [0.445, 0.461] from the
identical score vector. All six official labels collapse, by 0.205 to 0.307. The
patient-level permutation null on the same data is 0.523, so 0.456 is below what this
estimator reaches when there is by construction nothing to find. The result was
verified by four independent routes agreeing to 0.003, one of them an independent
reimplementation sharing no code, with its own folds, own seed, \texttt{sklearn}
AUROC and own bootstrap. Because the comparison is between two units of the same score
vector, this result does not depend on any published number being correct.

\textbf{2. Subject-clustered statistics, and the concrete cost of ignoring
clustering.} On the same data the naive slice-resampled interval is 1.5--2.0 times too
narrow; on simulated data with a closed-form true AUC, a nominal 95\% naive interval
covers the truth 46.5\% of the time against 91.5\% for the clustered one.

\textbf{3. A prevalence measurement with a defensible denominator.} A protocol-frozen
random-sample screen of a 9{,}979-record PubMed frame (SHA-256 \texttt{d611def0\ldots},
seed 20260729, both re-derivable offline), 250 records screened under a pre-specified
extension rule, four independent screeners on a fifteen-paper overlap set. Among 91
included and readable papers, \emph{no} paper reports a zero-image baseline with a
measured value. Because unreachability reached 32.6\% --- above the pre-specified 15\%
threshold --- \textbf{the protocol converts the headline from a point estimate to a
bounding interval of [0.0\%, 32.6\%]}, and we report it that way rather than promoting
the more flattering complete-case figure. The most robust form of the statement needs
no denominator at all: across 345 independently coded records covering 300 distinct
sampled papers, not one positive code appears on any zero-image sub-flag.

\textbf{4. A released tool.} \texttt{trivialbaselines}, CC~BY~4.0, depending on
\texttt{numpy} and \texttt{pandas} and nothing else. A benchmark can be audited by
someone who will never be granted its images, on a laptop, in one command.

\textbf{5. The formalisation and the slice-versus-patient contrast.} What the prior
literature measures is the cost of a \emph{wrong}, leaky split. Every number in this
paper was obtained under a \emph{correct}, patient-disjoint split. Stating that
distinction explicitly is necessary, because without it a reviewer would reasonably
reject the work as redundant with Yagis and Tampu.

\subsection*{What we do not claim}

We would rather declare the boundaries than have them found.

\textbf{Our pre-specified rank-inversion test is a null, and is reported as one.} Zero
of 447 sign-flipping method pairs survive on any of five cohorts. One cohort looked
spectacular --- Spearman $\rho = -0.31$ with 121 of 210 pairs flipping --- until the
noise floor was applied: split-half Kendall's $\tau$ at patient level is $-0.42$, so
two disjoint halves of subjects rank the methods worse than unrelated. The verdict is
\textsc{cannot rank}, and the near-miss is kept in the paper because it is the clearest
demonstration of why a noise floor is necessary. What does survive is narrower and is
stated narrowly: with architecture held fixed, aggregation to the patient unit shifts
the magnitude-versus-phase difference by $+0.028$ AUC [0.015, 0.041], $p = 0.001$,
positive for 6 of 6 architectures. \emph{The interaction is supported; the individual
sign flip is not.}

\textbf{Two of nine benchmark-arms do not fire, and are reported at equal
prominence.} LUNA16 returns a trivial fraction of $-0.002$, below its own random-score
reference, and PI-CAI's positional baseline lands on exactly 0.500. Two negatives out
of nine is an existence proof that the measure can fail to fire; it is not a
specificity estimate and we do not offer it as one.

\textbf{Every claim is about an evaluation protocol, never about what a model
learned.} We could not reproduce the pipeline of the benchmark we audit most closely
--- its protocol run on our data gives 0.616 against a reported 0.809 --- so we make
no claim about that implementation.

\textbf{The screen is not registered.} No OSF, PROSPERO or other deposit was made
before screening. The protocol, codebook, frame, permutation and sample were committed
to a public repository before any analysis-sample record was coded, which is a weaker
and different guarantee; any deposit made now is retrospective and is labelled so. We
therefore describe the screen as \emph{protocol-frozen} rather than
\emph{pre-registered} throughout, and use the latter word only with an explicit
qualifier. Supplementary Section~S7 gives the timeline, the digests, and a list of what
the git evidence does \emph{not} establish. Supplementary Section~S6 is a PRISMA 2020
checklist adapted for meta-research in which eight items are marked \emph{not done},
including four protocol deviations that no amount of rewriting closes.

\textbf{We make no regulatory claim.} A scoping sample of 30 cleared radiology AI
510(k) summaries is reported in Supplementary Section~S5 because we ran it, and its
result is negative for the hypothesis it was designed to test: the evaluation unit is
determinable in 28 of 30, no summary reports a slice-level metric, and the aggregation
step we recommend is already the regulatory requirement. We state that plainly rather
than omitting the analysis.

\subsection*{Fit with the journal}

The audience for this work is the reader who evaluates or reviews imaging AI. The
practical output is a seven-rule reporting protocol with a one-page checklist, each
rule written because a concrete failure was measured, and a tool that lets a reviewer
run the check on a benchmark's published label file in one command --- with no images,
no data-use agreement and no GPU. We believe that is squarely within \emph{Radiology:
Artificial Intelligence}'s remit.

\subsection*{Declarations}

The manuscript is original, is not under consideration elsewhere, and has not been
published in whole or in part. All four authors have approved the submission and meet
the ICMJE authorship criteria. No patient data were newly collected; every dataset used
is a public research release accessed under its own terms, and no unauthorised source
was used to obtain any full text at any stage. Funding and competing-interest
statements are supplied on the title page. All code, the frozen protocol and codebook,
the sampling frame with its digest, the seeded permutation, the sealed screener
submissions and every analysis output are released at
<<< REPOSITORY URL / ARCHIVAL DOI >>>.

We suggest that reviewers with expertise in \emph{neuroradiology and head CT triage},
\emph{meta-research and reporting standards}, and \emph{benchmark and challenge design}
would be best placed to assess this work. We have no reviewer exclusions to request.

Thank you for considering the manuscript.

\vspace{1.5em}

\noindent Yours sincerely,

\vspace{2.5em}

\noindent
\textbf{Sathvik Loke}\\
on behalf of Sathvik Loke, Ethan Johnson, Neeraj Movva and Aditya Raut

\end{document}
